%! AP Statistics — Unit 4, Lesson 4.9 — Guided Notes
\documentclass[10pt]{article}
\usepackage{apstats-article}
\usepackage{apstats-boxes}
\newcommand{\TallMath}[1]{$\displaystyle #1\rule[-1.4em]{0pt}{3.2em}$}

\begin{document}

\pageheader{Unit 4, Lesson 4.9}{Guided Notes}
\namedateperiod

\vspace{0.12in}

\begin{objectivebox}
\small By the end of this lesson, I will be able to\dots\\[4pt]
\writeline\\[1pt]
\writeline\\[1pt]
\writeline
\end{objectivebox}

\vspace{0.1in}

\begin{vocabbox}
\termblanklong{Linear transformation ($W = a + bX$)}
\termblanklong{Sum of random variables}
\termblanklong{Difference of random variables}
\termblanklong{Variance addition rule}
\end{vocabbox}

\vspace{0.1in}

\begin{hookbox}
\small
Maya's lunch order: $\mu_X = \$6.00$, $\sigma_X = \$1.50$. \quad
Jordan's lunch order: $\mu_Y = \$4.00$, $\sigma_Y = \$1.00$.\\[4pt]
\textbf{Expected total bill:} \blank{4cm} \quad
\textbf{SD of total bill:} \blank{4cm}\\[4pt]
\textit{(Try adding the SDs first. Then try the correct method. Are they the same?)}
\end{hookbox}

\vspace{0.1in}

\begin{notesbox}{Rule 1: Linear Transformation --- $W = a + bX$}
\small

Adding $a$ to $X$ \blank{3cm} the mean by $a$ but does \blank{3cm} change the spread.\\
Multiplying $X$ by $b$ \blank{3cm} both the mean and the SD by $b$.

\vspace{8pt}
\begin{center}
\renewcommand{\arraystretch}{1.6}
\begin{tabular}{@{}ll@{}}
\hline
$\mu_W = $ & \blank{6cm} \\
$\sigma_W = $ & \blank{6cm} \\
\hline
\end{tabular}
\end{center}

\vspace{8pt}
\textbf{Example:} $X$ = goals per game; $\mu_X = 0.90$, $\sigma_X = 0.70$.
Let $W = 3X$ (points per game, 3 points per goal).

$\mu_W = $ \blank{4cm} \quad $\sigma_W = $ \blank{4cm}

\end{notesbox}

\vspace{0.1in}

\begin{notesbox}{Rule 2: Sum of Independent Random Variables --- $T = X + Y$}
\small

\begin{center}
\renewcommand{\arraystretch}{1.6}
\begin{tabular}{@{}ll@{}}
\hline
$\mu_T = $ & \blank{6cm} \\
$\sigma^2_T = $ & \blank{8cm} (add \textbf{variances}, not SDs!) \\
$\sigma_T = $ & \blank{6cm} \\
\hline
\end{tabular}
\end{center}

\vspace{8pt}
\textbf{Example:} Maya's bill ($\mu_X=6$, $\sigma^2_X=2.25$) + Jordan's bill ($\mu_Y=4$, $\sigma^2_Y=1.00$).

$\mu_T = $ \blank{3cm} \quad $\sigma^2_T = $ \blank{3cm} \quad $\sigma_T = \sqrt{\blank{2cm}} = $ \blank{3cm}

\end{notesbox}

\vspace{0.1in}

\begin{notesbox}{Rule 3: Difference of Independent Random Variables --- $D = X - Y$}
\small

\begin{center}
\renewcommand{\arraystretch}{1.6}
\begin{tabular}{@{}ll@{}}
\hline
$\mu_D = $ & \blank{6cm} \\
$\sigma^2_D = $ & \blank{8cm} (variances \textbf{still add}!) \\
$\sigma_D = $ & \blank{6cm} \\
\hline
\end{tabular}
\end{center}

\vspace{8pt}
\textbf{Key insight:} When you subtract two random variables, the variance
\blank{3cm} (increases / decreases / stays the same).

\vspace{8pt}
\textbf{Example:} Difference in bills (Maya $-$ Jordan).

$\mu_D = $ \blank{3cm} \quad $\sigma^2_D = $ \blank{3cm} \quad $\sigma_D = $ \blank{3cm}

Compare: $\sigma_T = $ \blank{2cm} and $\sigma_D = $ \blank{2cm}. Are they equal?  \blank{2cm}

\end{notesbox}

\vspace{0.1in}

\begin{practicebox}
\small
$X$: $\mu_X = 50$, $\sigma_X = 8$, $\sigma^2_X = 64$.
$Y$: $\mu_Y = 30$, $\sigma_Y = 5$, $\sigma^2_Y = 25$. ($X$ and $Y$ are independent.)

\begin{enumerate}[leftmargin=1.5em, itemsep=8pt]
  \item $T = X + Y$: \quad $\mu_T = $ \blank{2cm} \quad $\sigma^2_T = $ \blank{2cm} \quad
        $\sigma_T = $ \blank{3cm}

  \item $D = X - Y$: \quad $\mu_D = $ \blank{2cm} \quad $\sigma^2_D = $ \blank{2cm} \quad
        $\sigma_D = $ \blank{3cm}

  \item $W = 10 + 2X$: \quad $\mu_W = $ \blank{3cm} \quad $\sigma_W = $ \blank{3cm}
\end{enumerate}
\end{practicebox}

\end{document}
